\begin{lemma}
Let \(\Lambda\) be a finite label set and let \(B\) be a
null-measurable base event in a measured space.  For each
\(\lambda\in\Lambda\), let \(Q_\lambda\) be a null-measurable label event.
For a truth vector \(\tau\subseteq\Lambda\), define
\[
  A_\tau
  =
  B\cap
  \{s:s\in Q_\lambda\Longleftrightarrow \lambda\in\tau
       \text{ for every }\lambda\in\Lambda\}.
\]
Then the finite family \((A_\tau)_{\tau\subseteq\Lambda}\) covers \(B\), is
pairwise almost-disjoint, and consists of null-measurable sets.
\end{lemma}
\begin{proof}
For each point \(s\in B\), define its truth vector by
\[
  \tau(s)=\{\lambda\in\Lambda:s\in Q_\lambda\}.
\]
Since \(\Lambda\) is finite, \(\tau(s)\) is one of the allowed finite truth
vectors.  By construction, \(s\in A_{\tau(s)}\).  This proves
\[
  B\subseteq \bigcup_{\tau\subseteq\Lambda} A_\tau .
\]
The reverse inclusion is immediate from the definition \(A_\tau\subseteq B\).
Thus the atoms cover \(B\).

If \(\tau_1\ne\tau_2\), then some label
\(\lambda\in\Lambda\) belongs to exactly one of the two truth vectors.  A
point in \(A_{\tau_1}\cap A_{\tau_2}\) would have to satisfy both
\[
  s\in Q_\lambda\Longleftrightarrow \lambda\in\tau_1
  \quad\text{and}\quad
  s\in Q_\lambda\Longleftrightarrow \lambda\in\tau_2,
\]
which is impossible.  Hence \(A_{\tau_1}\cap A_{\tau_2}=\emptyset\) for
distinct truth vectors.  Empty intersection implies almost-disjointness for
every measure.

Finally,
\[
  \{s:s\in Q_\lambda\Longleftrightarrow \lambda\in\tau
       \text{ for every }\lambda\}
  =
  \bigcap_{\lambda\in\Lambda}
  \begin{cases}
    Q_\lambda, & \lambda\in\tau,\\
    Q_\lambda^c, & \lambda\notin\tau .
  \end{cases}
\]
This is a finite intersection of null-measurable sets and complements of
null-measurable sets, so it is null-measurable.  Intersecting it with the
null-measurable base event \(B\) proves that each \(A_\tau\) is
null-measurable.
\end{proof}
